\documentclass[11pt,a4paper]{article}

\usepackage[margin=2.5cm]{geometry}
\usepackage{amsmath,amssymb,amsthm}
\usepackage{booktabs}
\usepackage[hidelinks]{hyperref}

\title{Stability of Scenario-Based Stochastic Facility Location:\\
Scenario Generation and Stabilization}
\author{Akeem Alabi Akinyemi}
\date{Work in progress}

\begin{document}
\maketitle

\begin{abstract}
We study a two-stage stochastic facility location problem in which sites are opened
before demand is known and are chosen using a finite set of demand scenarios.
This note will report how stable the chosen sites and their cost are when the
scenarios change, how scenario generation affects the result, and whether
stabilization techniques help. This version contains the problem statement and
the model only.
\end{abstract}

\section{Problem statement}

Facility location under uncertain demand can be modelled as a two-stage stochastic
mixed-integer program, in which site-opening decisions are made before demand is
revealed and are approximated using a finite set of demand scenarios. In practice
the chosen sites depend strongly on which scenarios are sampled, and plans optimized
on one sample tend to perform worse on new demand. We investigate three questions:
\begin{enumerate}
  \item How unstable are the optimal site selection and cost as the number and
        sampling of scenarios change, and how large is the gap between promised and
        realized cost?
  \item How does the method used to generate demand scenarios (independent, spatially
        correlated, resampled from data, and reduced by clustering) affect stability
        and out-of-sample cost?
  \item Can stabilization techniques (consensus over resampled solutions, a
        risk-averse CVaR objective, and a regularized version) reduce instability and
        the out-of-sample gap, and at what cost in expected total cost?
\end{enumerate}
Results are compared against a mean-demand deterministic plan, a simple heuristic,
and a large-sample reference solution.

\section{The two-stage model}

\subsection{Sets and data}

Let $I$ be the set of candidate sites and $J$ the set of demand zones. A finite set
$S$ of $N$ demand scenarios is given, each with probability $1/N$. The data are
\begin{center}
\begin{tabular}{@{}ll@{}}
\toprule
Symbol & Meaning \\
\midrule
$f_i$    & cost of opening site $i$ \\
$K_i$    & capacity of site $i$ \\
$c_{ij}$ & cost of serving one unit of demand from site $i$ to zone $j$ \\
$\pi$    & penalty per unit of unmet demand \\
$d_{js}$ & demand of zone $j$ in scenario $s$ \\
\bottomrule
\end{tabular}
\end{center}
We assume $\pi > \max_{i,j} c_{ij}$, so that serving a unit is always cheaper than
leaving it unmet when capacity is available.

\subsection{Decision variables}

The first-stage variables are
\[
  y_i =
  \begin{cases}
    1 & \text{if site } i \text{ is opened},\\
    0 & \text{otherwise,}
  \end{cases}
  \qquad i \in I.
\]
They carry no scenario index, because sites are opened once, before demand is known.
The second-stage variables are $x_{ijs} \ge 0$, the number of units served from site
$i$ to zone $j$ in scenario $s$, and $u_{js} \ge 0$, the number of units of unmet
demand in zone $j$ in scenario $s$. These depend on the scenario because they are
chosen after demand is revealed.

\subsection{Formulation}

The sample average approximation of the problem is
\begin{align}
  \min_{y,\,x,\,u} \quad
    & \sum_{i \in I} f_i\, y_i
      + \frac{1}{N} \sum_{s \in S} \Bigl[ \sum_{i \in I} \sum_{j \in J} c_{ij}\, x_{ijs}
      + \pi \sum_{j \in J} u_{js} \Bigr] \label{eq:obj}\\[4pt]
  \text{s.t.} \quad
    & \sum_{i \in I} x_{ijs} + u_{js} = d_{js}
      && \forall j \in J,\ s \in S, \label{eq:balance}\\
    & \sum_{j \in J} x_{ijs} \le K_i\, y_i
      && \forall i \in I,\ s \in S, \label{eq:capacity}\\
    & y_i \in \{0,1\},\quad x_{ijs} \ge 0,\quad u_{js} \ge 0. \label{eq:domain}
\end{align}

The first term of \eqref{eq:obj} is the fixed cost of the opened sites, paid once.
The second term is the average, over scenarios, of the serving cost and the penalty
for unmet demand. Constraint \eqref{eq:balance} requires every unit of demand in a
zone to be either served or counted as unmet. Constraint \eqref{eq:capacity} limits
what a site can serve to its capacity, and to zero if it is not opened; it is the
only constraint linking the first stage to the second.

\subsection{Properties}

\paragraph{Complete recourse.}
For every choice of $y$ and every scenario, setting $x = 0$ and $u_{js} = d_{js}$
satisfies \eqref{eq:balance} and \eqref{eq:capacity}. The second-stage problem is
therefore always feasible, and the model is never infeasible.

\paragraph{Separability.}
For fixed $y$, the scenarios do not interact: the problem decomposes into $N$
independent linear programs, one per scenario. Writing $Q(y,s)$ for the optimal value of
scenario $s$, the objective becomes
\[
  \sum_{i \in I} f_i\, y_i + \frac{1}{N} \sum_{s \in S} Q(y,s).
\]
This allows a plan $y$ to be evaluated on a large test set by solving one small
linear program per test scenario.

\paragraph{Size.}
With $|I|$ sites, $|J|$ zones and $N$ scenarios, the model has $|I|$ binary variables,
$N|I||J| + N|J|$ continuous variables, and $N(|J| + |I|)$ constraints
(excluding bounds). For $|I|=12$, $|J|=40$ and $N=50$ this is $12$ binary variables,
$26{,}000$ continuous variables and $2{,}600$ constraints.

\end{document}